Sigmalized residual calibrations for this analysis were performed for EMCal(e,w)(0-3) in Run14 \HeAu, EMCal(e,w)(0-3), TOFe, and TOFw in Run15 \pAl, TOFe and TOFw in Run15 \pAu.

Matching cuts are important in this analysis for improving the signal to background ratio.
The sigmalized residuals that are used in these cut are defined as follows

\begin{equation}
   sdz = \frac{(dz - \mu_{dz})}{\sigma_{dz}}, \ \ \ sd\varphi = \frac{(d\varphi - \mu_{d\varphi})}{\sigma_{d\varphi}}
\end{equation}

Where $dz$ and $d\varphi$ are deviation of detected $z$ and $\varphi$ positions respectively from the predicted by the projection of a DC-PC1 track on the detector plane,
$\mu_{dz}$, $\mu_{d\varphi}$, $\sigma_{dz}$, $\sigma_{d\varphi}$ are means and sigmas of $dz$ and $dphi$ distributions.
The calibration of sigmalized residuals is a process of parametrization of these means and sigmas in terms of centrality, $z_{\text{DC}}$, and \pT. 

The calibration was performed by dissecting the dataset in 2 charge bins (positive and negative), 10 $z_{\text{DC}}$ bins (Table. \ref{table:sigm_res_cal_zdc_bins}) for all datasets, and variable number and range of centrality bins depending on the dataset: 17 bins for Run14 \HeAu (Table. \ref{table:sigm_res_cal_centr_bins_run14heau}), 7 bins for Run15 \pAl (Table. \ref{table:sigm_res_cal_centr_bins_run15pal}), 8 bins for Run15 \pAu (Table. \ref{table:sigm_res_cal_centr_bins_run15pau}).

\begin{table}[h!]
	\centering
	\begin{tabular}{c||c|c|c|c|c|c|c|c|c|c}
      Bin minimum & -75 & -60 & -45 & -30 & -15 & 0 & 15 & 30 & 45 & 60 \\
      \hline
      Bin maximum & -60 & -45 & -30 & -15 & 0 & 15 & 30 & 45 & 60 & 75 \\
	\end{tabular}
   \caption{$z_{\text{DC}}$ bins for sigmalized residuals calibration}
	\label{table:sigm_res_cal_zdc_bins}
\end{table}

\begin{table}[h!]
	\centering
	\begin{tabular}{c||c|c|c|c|c|c|c|c|c|c|c|c|c|c|c|c|c}
      Bin minimum & 0 & 5 & 10 & 15 & 20 & 25 & 30 & 35 & 40 & 45 & 50 & 55 & 60 & 65 & 70 & 75 & 80 \\
      \hline
      Bin maximum & 5 & 10 & 15 & 20 & 25 & 30 & 35 & 40 & 45 & 50 & 55 & 60 & 65 & 70 & 75 & 80 & 88 \\
	\end{tabular}
   \caption{Centrality bins for sigmalized residuals calibration in Run14 \HeAu}
	\label{table:sigm_res_cal_centr_bins_run14heau}
\end{table}

\begin{table}[h!]
	\centering
	\begin{tabular}{c||c|c|c|c|c|c|c}
      Bin minimum & 0 & 10 & 20 & 30 & 40 & 50 & 60 \\
      \hline
      Bin maximum & 10 & 20 & 30 & 40 & 50 & 60 & 72 \\
	\end{tabular}
   \caption{Centrality bins for sigmalized residuals calibration in Run15 \pAl}
	\label{table:sigm_res_cal_centr_bins_run15pal}
\end{table}

\begin{table}[h!]
	\centering
	\begin{tabular}{c||c|c|c|c|c|c|c|c}
      Bin minimum & 0 & 10 & 20 & 30 & 40 & 50 & 60 & 70 \\
      \hline
      Bin maximum & 10 & 20 & 30 & 40 & 50 & 60 & 70 & 84 \\
	\end{tabular}
   \caption{Centrality bins for sigmalized residuals calibration in Run15 \pAu}
	\label{table:sigm_res_cal_centr_bins_run15pau}
\end{table}

Additionally each centrality and $z_{\text{DC}}$ bins were dissected into fine $p_T$ bins with widths 0.1 GeV for $p_T \lesssim 1.5$ GeV, 0.2 GeV for $1.5 \lesssim p_T 2 \div 2.5$ and 0.5 GeV for $2 \div 2.5 \lesssim p_T \lesssim 3\div3.5$ depending on the statistics. 
Higher $p_T$ regions were not used for calibration due to high background contamination and low statistics which leads to unpredictable behavior of approximation algorithms. 
In each charge-centrality-$z_{\text{DC}}$-$p_T$ bin 1-D distributions of $dz$ and $d\varphi$ were approximated with the sum of 2 gausses - one for the signal and the other for the background. 
Example of these approximations are presented on Fig. \ref{fig:sim_res_cal_dz_fit_example}. 
Means and sigmas for different $p_T$ were then approximated separately in each charge-centrality-$z_{\text{DC}}$ bin.

\begin{figure}[h!]
   \centering
   \includegraphics[width=\textwidth]{../../../CalPHENIX/output/SigmalizedResiduals/Run15pAl200/EMCale2/dz_pos_c0-10_zDC0-15.pdf}
   \caption[Example of $dz$ approximations]
   {Example of EMCale2 $dz$ approximations in Run15pAl for 0-10\% centrality class for $0 < z_{\text{DC}} < 15$, $\text{charge} > 0$}
   \label{fig:sim_res_cal_dz_fit_example}
\end{figure}

It was found that same fit function forms for approximation means and sigmas of $dz$ and $d\varphi$ work universally well for different datasets, so the same fit types were used for the same detector across different datasets. 
Fit functions for sigmalized residuals parametrization for different detectors are presented in Table. \ref{table:sigm_res_cal_fit_func}. Example of $\mu_{dz}$ approximations are shown on Fig. \ref{fig:sim_res_cal_means_dz_fit_example}.

\begin{table}[h!]
	\centering
	\begin{tabular}{c||c|c}
       & TOFe \& TOFw \& EMCale(0-1) & EMCale(2-3) \& EMCalw(0-3) \\
      \hline
      \hline
      $\mu_{d\varphi}$ & $p_0 + \frac{p_1}{p_T} + \frac{p_2}{p_T^2} + \frac{p_3}{p_T^3}$ & $p_0 + \frac{p_1}{p_T^2} + \frac{p_2}{p_T^3} + e^{p_3 + p_4 p_T} + p_5 p_T$ \\
      \hline
      $\mu_{dz}$ & $p_0 + \frac{p_1}{p_T^2} + p_2 p_T$ & $p_0 + \frac{p_1}{p_T} + p_2 p_T$ \\
      \hline
      $\sigma_{d\varphi}$ & $p_0 + \frac{p_1}{p_T^2} + p_2 p_T$ & $p_0 + \frac{p_1}{p_T} + p_2 p_T$ \\
      \hline
      $\sigma_{dz}$ & $p_0 + \frac{p_1}{p_T^2} + p_2 p_T$ & $p_0 + \frac{p_1}{p_T} + p_2 p_T$ \\
	\end{tabular}
   \caption[$\mu_{dz}$, $\mu_{d\varphi}$, $\sigma_{dz}$, and $\sigma_{d\varphi}$ parametrization]
           {$\mu_{dz}$, $\mu_{d\varphi}$, $\sigma_{dz}$, and $\sigma_{d\varphi}$ parametrization as functions of \pT for different detectors}
	\label{table:sigm_res_cal_fit_func}
\end{table}

\begin{figure}[h!]
   \centering
   \begin{subfigure}{0.48\textwidth}
      \includegraphics[width=\textwidth]{../../../CalPHENIX/output/SigmalizedResiduals/Run15pAl200/EMCale2/means_dz_pos_c0-10.pdf}
      \subcaption{}
   \end{subfigure}
   \hfill
   \begin{subfigure}{0.48\textwidth}
      \includegraphics[width=\textwidth]{../../../CalPHENIX/output/SigmalizedResiduals/Run15pAl200/EMCale2/sigmas_dz_pos_c0-10.pdf}
      \subcaption{}
   \end{subfigure}
   \caption[Example of $\mu_{dz}$ and $\sigma_{dz}$ approximations]
           {Example of EMCale2 $\mu_{dz}$ (a) and $\sigma_{dz}$ (b) approximations in Run15 \pAl in 0-10\% centrality class for different $z_{\text{DC}}$ bins, $\text{charge} > 0$}
   \label{fig:sim_res_cal_means_dz_fit_example}
\end{figure}

While calibrations were performed in $0.3 < p_T < 2\div3$ range, parametrizations can be used on higher \pT due to the trend of means and sigmas of $dz$ and $d\varphi$ becoming constants at $p_T \gtrsim 2\div2.5$. 
Thus the parametrization of means and sigmas at the highest $p_T$ measured during the calibration was used to estimate the same parameters at higher $p_T$.

After obtaining all $\mu_{dz}$, $\mu_{d\varphi}$, $\sigma_{dz}$, and $\sigma_{d\varphi}$ parametrizations, $sdz$, $sd\varphi$ distributions were investigated for possible additional recalibrations in each charge-centrality-$z_{\text{DC}}$-$p_T$. 
Since these distributions are sigmalized, the main signal should be a gauss centered around 0 with width 1.
However approximations of $\mu_{dz}$, $\mu_{d\varphi}$, $\sigma_{dz}$, and $\sigma_{d\varphi}$ were not ideal and for some parametrizations additional small constant shift for means and small constant scaling for sigmas were implemented since non-constant deviations were negligible or within uncertainties. Fig. \ref{fig:sim_res_recal_means_sdz_fit_example} shows an example of these constant offsets.

\begin{figure}[h!]
   \centering
   \begin{subfigure}{0.48\textwidth}
      \includegraphics[width=\textwidth]{../../../CalPHENIX/output/SigmalizedResiduals/Run14HeAu200/EMCale0/means_sdz_neg_c15-20.pdf}
      \subcaption{}
   \end{subfigure}
   \hfill
   \begin{subfigure}{0.48\textwidth}
      \includegraphics[width=\textwidth]{../../../CalPHENIX/output/SigmalizedResiduals/Run14HeAu200/EMCale0/sigmas_sdz_neg_c15-20.pdf}
      \subcaption{}
   \end{subfigure}
   \caption[Example of $\mu_{sdz}$ and $\sigma_{sdz}$ constant offsets]
           {Example of EMCale0 $\mu_{sdz}$ (a) and $\sigma_{sdz}$ (b) approximations in Run14 \HeAu in 15-20\% centrality class for different $z_{\text{DC}}$ bins, $\text{charge} < 0$. Constant offset can fix the deviation of $\mu_{sdz}$ from 0, and constant scaling can fix $\sigma_{sdz}$ deviation from 1.}
   \label{fig:sim_res_recal_means_sdz_fit_example}
\end{figure}

An example of applying sigmalized residuals calibrations is shown on Fig. \ref{fig:sdz_example}.

\begin{figure}[h!]
   \centering
   \begin{subfigure}{0.48\textwidth}
      \includegraphics[width=\textwidth]{other/sdphi_example.pdf}
      \subcaption{}
   \end{subfigure}
   \hfill
   \begin{subfigure}{0.48\textwidth}
      \includegraphics[width=\textwidth]{other/sdz_example.pdf}
      \subcaption{}
   \end{subfigure}
   \caption[Example of $sd\varphi$ and $sdz$ distributions after applying calibrations]
           {Example of EMCale0 $sd\varphi$ (a) and $sdz$ (b) distributions after applying sigmalized residuals calibrations in Run15 \pAl, $\text{charge} > 0$}
   \label{fig:sim_res_recal_means_sdz_fit_example}
\end{figure}
